\documentclass[11pt,a4paper]{article}

\usepackage[margin=25mm]{geometry}
\usepackage[T1]{fontenc}
\usepackage{lmodern}
\usepackage{amsmath,amssymb,amsthm}
\newcommand{\Area}{\operatorname{area}}
\theoremstyle{definition}
\newtheorem{definition}{Definition}
\newtheorem{problem}{Problem}

\setlength{\parindent}{0pt}
\setlength{\parskip}{6pt}

\begin{document}

\begin{center}
  {\Large\bfseries Heilbronn's Triangle Problem: The Upper Bound}
\end{center}
\section{Definition}
Let $U=[0,1]^2$. For a finite set $P\subseteq U$, define
\begin{equation}\label{eq:delta}
 h(P)
 =\min_{\substack{x,y,z\in P\\\text{distinct}}}
   \Area\bigl(\Delta(x,y,z)\bigr),
\end{equation}
where $\Delta(x,y,z)$ denotes the triangle formed by the three points $x,y,z$,
and degenerate triangles have area zero.

\begin{definition}[The $n$th Heilbronn number]
For $n\geq 3$, the \emph{$n$th Heilbronn number} of the unit square is
\begin{equation}\label{eq:H}
 H(n)=\max_{\substack{P\subseteq U\\|P|=n}} h(P).
\end{equation}
A set attaining this maximum is called an \emph{optimal Heilbronn
configuration}.
\end{definition}

The maximum exists by compactness. The use of the unit square is convenient but
not essential for asymptotic questions: for any two fixed convex bodies with
nonempty interior, the corresponding Heilbronn numbers differ by multiplicative
constants depending only on the two bodies.


\section{Current best result}
The best known asymptotic bounds are
\begin{equation}
  H(n)=\Omega\!\left(\frac{\log n}{n^2}\right),
  \qquad
  H(n)\leq n^{-7/6+o(1)}.
\end{equation}


\section{Problem formalization}

Prove an asymptotic upper bound for $H(n)$ that is strictly stronger than
\begin{equation}
  H(n)\leq n^{-7/6+o(1)}.
\end{equation}

\end{document}
